\chapter{Simulation Study}
\label{ch:simulation}

This chapter exercises the framework in closed loop on the full nonlinear plants
of \cref{ch:hexacopter,ch:fixedwing}. It reports five studies: a representative
six-hexacopter mission; a Monte-Carlo robustness campaign over noise
realisations, GNSS-denial severity, degraded communication graphs and an
estimator ablation, including the multi-hop anchoring extension of
\cref{sec:dmhe-fusion}; a computational study of the embedded per-agent solvers
instantiated on the drone models; a closed-loop validation with those embedded
solvers \emph{in} the loop; and a four-aircraft fixed-wing formation.
Every number is produced by the accompanying code; the controllers and
estimators see only noisy measurements.

\section{Hexacopter cluster: representative mission}
\label{sec:sim-hexa}

\paragraph{Setup.} Six hexacopters (parameters in \cref{tab:hexa-params}) start
scattered on the ground and are commanded through four mission phases: vertical
take-off, assembly into a ring formation of radius \SI{3}{\meter}, a sinusoidal
translation of the formation centre, and a mid-mission reconfiguration in which
the ring contracts and rotates. The communication graph is a ring augmented with
second-neighbour chords (degree four, algebraic connectivity
$\lambda_2(\Lap)=4.0$). Each agent runs the full stack of
\cref{alg:dmhe,alg:dmpc}: distributed estimation from noisy GNSS
($\sigma=\SI{0.3}{\meter}$), inertial measurement and relative measurements
($\sigma=\SI{0.1}{\meter}$); a consensus on the formation centre; a
model-predictive position command realised through differential flatness and a
geometric inner attitude loop on the $13$-state quaternion plant; and the CBF
safety filter of \cref{sec:cbf} with safety radius $D=\SI{1.4}{\meter}$. To test
the DMHE, agents $3,4,5$ suffer a complete GNSS blackout during $t\in[12,22]\,$s
and must localise by relative fusion to their still-tracked neighbours.

\paragraph{Formation and safety.} \Cref{fig:hexa-traj} shows the closed-loop
trajectories of a representative run: the cluster takes off, assembles the ring,
follows the swept centre and contracts into the reconfigured formation.
\Cref{fig:hexa-formation} quantifies the behaviour. The formation-\emph{shape}
error (deviation from the ring geometry about the achieved centroid) settles
near \SI{0.09}{\meter}; the total formation error, which also carries the
uniform tracking lag of the moving centre, oscillates with the sweep and stays
below \SI{4}{\meter}. The minimum pairwise separation remains above the
\SI{1.4}{\meter} safety radius for the entire mission, including the
contraction, confirming that the safety filter keeps the transient
collision-free.

\begin{figure}[t]
\centering
\includegraphics[width=0.82\linewidth]{figs/fig_hexa_traj3d.pdf}
\caption{Closed-loop trajectories of the six-hexacopter cluster on the full
nonlinear quaternion plants (circles: start; triangles: end). The mission
comprises take-off, ring assembly, a swept translation and a contract-and-rotate
reconfiguration.}
\label{fig:hexa-traj}
\end{figure}

\begin{figure}[t]
\centering
\includegraphics[width=0.86\linewidth]{figs/fig_hexa_formation.pdf}
\caption{Formation tracking error (top) and minimum inter-agent separation
(bottom) for the representative run. The separation stays above the safety
radius $D=\SI{1.4}{\meter}$ throughout; the shaded band marks the GNSS blackout
of agents $3,4,5$.}
\label{fig:hexa-formation}
\end{figure}

\paragraph{Distributed estimation through the blackout.} \Cref{fig:hexa-dmhe}
shows the estimation result. During the ten-second blackout the local
filter---which can only dead-reckon on inertial measurements once its GNSS fix
is lost---drifts steadily past \SI{1.2}{\meter}, while the distributed
estimator, fusing relative measurements to the three neighbours that retain
GNSS, holds the denied agents' position error essentially flat at the
relative-measurement noise level. The accurate estimate is what keeps the
control loop stable and the formation intact through the outage; the closed-loop
consequence of \emph{not} having it is quantified by the ablation in
\cref{sec:sim-mc}. \Cref{fig:hexa-timing} shows the per-step wall time of the
reference-layer estimation and coordination for the whole six-agent cluster on
one core (median \SI{\sim0.35}{\milli\second}); the embedded per-agent solves
are measured separately in \cref{sec:sim-compute}.

\begin{figure}[t]
\centering
\includegraphics[width=0.82\linewidth]{figs/fig_hexa_dmhe.pdf}
\caption{Position RMSE of the GNSS-denied agents (representative run). Outside
the blackout the distributed estimator and the local filter coincide; inside it
(shaded) the local filter dead-reckons and drifts, while the DMHE holds the
error at the relative-measurement noise level by anchoring to GNSS-tracked
neighbours.}
\label{fig:hexa-dmhe}
\end{figure}

\begin{figure}[t]
\centering
\includegraphics[width=0.74\linewidth]{figs/fig_hexa_timing.pdf}
\caption{Per-control-step wall time of the reference-layer distributed
estimation and coordination for the full six-agent cluster on one core.}
\label{fig:hexa-timing}
\end{figure}

\section{Monte-Carlo robustness study}
\label{sec:sim-mc}

A single trajectory establishes feasibility, not robustness. The full mission of
\cref{sec:sim-hexa} was therefore repeated in a $126$-run Monte-Carlo campaign:
$30$ independent noise realisations of the nominal scenario, a GNSS-denial
severity sweep ($12$ seeds for each contiguous denied set of size $1$--$5$ of
$6$), two communication-graph perturbations ($12$ seeds each), and a closed-loop
estimator ablation ($12$ seeds). Every run regenerates all measurement noise,
initial scatter and reference-sensing noise. \Cref{tab:mc} summarises the
campaign; across all $126$ runs there were \emph{no} divergences and \emph{no}
solver failures.

\begin{table}[t]
\centering
\caption{Monte-Carlo robustness campaign ($150$ runs; $10$-s GNSS blackout in
all scenarios). RMSE columns are the denied agents' position RMSE during the
blackout. ``Shape'' is the formation-shape error about the achieved centroid
over the assembled phase. ``Margin breach'' is time with minimum separation
below the $D=\SI{1.4}{\meter}$ design radius, per run (physical contact never
occurred; the smallest separation observed anywhere was \SI{0.75}{\meter}).}
\label{tab:mc}
\small
\begin{tabular}{@{}lccccccc@{}}
\toprule
& & \multicolumn{2}{c}{DMHE RMSE [m]} & Local filter & Shape & Min sep. & Margin\\
\cmidrule(lr){3-4}
Scenario & $n$ & mean\,/\,med & p95 & mean [m] & [m] & [m] & breach [s]\\
\midrule
Nominal ($3/6$ denied) & 30 & 0.091\,/\,0.091 & 0.099 & 0.96 & 0.094 & $\ge$1.55 & 0\\
Denied $1/6$ & 12 & 0.068\,/\,0.068 & 0.074 & 0.97 & 0.101 & $\ge$1.49 & 0\\
Denied $2/6$ & 12 & 0.078\,/\,0.078 & 0.085 & 0.90 & 0.094 & $\ge$1.54 & 0\\
Denied $4/6$ & 12 & 0.106\,/\,0.105 & 0.117 & 0.91 & 0.090 & $\ge$1.54 & 0\\
Denied $5/6$$^{\dagger}$ & 12 & 0.451\,/\,0.365 & 0.689 & 0.96 & 0.267 & $\ge$0.90 & 0.50\\
Ring graph$^{\dagger}$ ($\lambda_2{=}1$) & 12 & 0.506\,/\,0.474 & 0.799 & 0.95 & 0.244 & $\ge$0.98 & 0.20\\
Denied $5/6$, multi-hop$^{\S}$ & 12 & 0.119\,/\,0.120 & 0.126 & 0.98 & 0.069 & $\ge$1.49 & 0\\
Ring graph, multi-hop$^{\S}$ & 12 & 0.111\,/\,0.112 & 0.119 & 0.94 & 0.100 & $\ge$1.49 & 0\\
Edge loss $(3,2)$ in blackout & 12 & 0.094\,/\,0.092 & 0.107 & 0.89 & 0.095 & $\ge$1.50 & 0\\
Controller on local filter$^{\ddagger}$ & 12 & 0.089\,/\,0.087 & 0.099 & 0.91 & 0.449 & $\ge$0.75 & 0.80\\
\midrule
\multicolumn{8}{@{}l}{\footnotesize Divergences: $0/150$.\quad Safety-QP failures: $0$.\quad
Reference-layer compute: median \SI{0.35}{\milli\second}\,/\,step (6 agents).}\\
\multicolumn{8}{@{}p{0.94\linewidth}@{}}{\footnotesize
$^{\dagger}$Anchor-coverage violated: at $5/6$ denied, agent $3$'s neighbours
are all denied; on the plain ring, denied agent $4$'s neighbours are both
denied. The affected agent must dead-reckon (\cref{sec:sim-mc-graph}).\quad
$^{\S}$Same scenarios with the leveled multi-hop anchoring of
\cref{sec:dmhe-fusion} ($L{=}3$): coverage is restored and the margin breaches
vanish.\quad
$^{\ddagger}$Ablation: the DMHE still runs (its passive RMSE is the first
column) but the controller is fed the dead-reckoning local filter.}\\
\bottomrule
\end{tabular}
\end{table}

\paragraph{Noise robustness.} Over the $30$ nominal seeds the blackout
performance is remarkably tight: denied-agent DMHE RMSE of
$0.091$\,/\,$0.091$\,/\,$0.099$\,m (mean\,/\,median\,/\,95th percentile)
against $0.96$\,m for dead reckoning, formation-shape error
\SI{0.094}{\meter}, and a worst-case-over-all-seeds minimum separation of
\SI{1.55}{\meter}---the safety margin was never touched. The single
representative run of \cref{sec:sim-hexa} is thus typical, not favourable.

\paragraph{GNSS-denial severity.} \Cref{fig:denial-sweep} sweeps the denied set
from one to five of six agents. The DMHE error grows only mildly with severity
($0.068\to0.106$\,m through $4/6$ denied---more denied agents means fewer
anchors shared among more consumers) while dead reckoning sits an order of
magnitude higher throughout. At $5/6$ denied the one-hop curve breaks upward, for a
structural reason the campaign makes precise: denied agent $3$ then has no
GNSS-good agent within one hop, the one-hop fusion policy correctly refuses to
anchor it to other drifting agents, and it reverts to dead reckoning.
Resilience is thus governed not by the denial count as such but by
\emph{anchor coverage}. The leveled multi-hop anchoring of
\cref{sec:dmhe-fusion} repairs exactly this case: agent $3$ anchors, at level
two, to neighbours that are themselves level-one anchored to the single
remaining GNSS-good agent, and over $12$ seeds the denied-agent RMSE returns to
$0.119$\,m mean\,/\,$0.126$\,m p95 (third curve of \cref{fig:denial-sweep})
with the separation margin fully restored ($\ge$\SI{1.49}{\meter}, zero
breaches). Through $4/6$ denial the multi-hop and one-hop policies coincide by
construction, so the extension costs nothing when coverage already holds.

\begin{figure}[t]
\centering
\includegraphics[width=0.8\linewidth]{figs/fig_denial_sweep.pdf}
\caption{GNSS-denial severity sweep ($12$ Monte-Carlo seeds per point; bands are
min--max). One-hop DMHE degrades gracefully while anchor coverage holds and
breaks upward exactly when it fails (annotated); the leveled multi-hop
anchoring of \cref{sec:dmhe-fusion} restores coverage and keeps the error at
the decimetre level through $5/6$ denial; the dead-reckoning baseline is an
order of magnitude worse throughout.}
\label{fig:denial-sweep}
\end{figure}

\paragraph{Communication-graph robustness.}
\label{sec:sim-mc-graph}
Two graph perturbations probe the same condition from the topology side. Losing
the single edge $(3,2)$ for four seconds in the middle of the blackout is
absorbed almost without trace ($0.091\to0.094$\,m mean RMSE): agent $3$'s second
anchor covers the lost one, and the safety margin is untouched. Replacing the
chorded graph by the plain ring ($\lambda_2$ drops from $4$ to $1$) is not:
denied agent $4$'s ring neighbours are themselves both denied, anchor coverage
fails, its estimate drifts, and the eroded estimate quality lets the true
separation dip below the design margin for $\sim$\SI{0.2}{\second} per run
(minimum \SI{0.98}{\meter}; no contact). The two experiments bracket the design
rule: under one-hop fusion the graph must be provisioned so that plausible
denial sets keep every agent within one hop of an anchor---degree, not only
algebraic connectivity, carries the estimation-side robustness. Multi-hop
anchoring relaxes the rule to reachability: on the same ring, denied agent $4$
anchors at level two through its level-one neighbours, and over $12$ seeds the
ring scenario recovers to $0.111$\,m mean\,/\,$0.119$\,m p95 with the margin
restored ($\ge$\SI{1.49}{\meter}, zero breaches; \cref{tab:mc}). A sparse graph
then costs accuracy only through the longer anchor chains, not through outright
coverage failure.

\paragraph{Closed-loop value of the DMHE.} The ablation feeds the controller the
dead-reckoning local filter while the DMHE runs passively. The consequence of
the $\sim$\SI{0.9}{\meter} estimate drift is now dynamical, not merely
statistical: the formation-shape error during the blackout grows to
\SI{0.449}{\meter} ($4.8\times$ the nominal \SI{0.094}{\meter}), the peak
formation error rises to $4.2$\,m mean\,/\,$4.9$\,m p95 (nominal:
$3.9$\,/\,$4.0$), and---because the safety filter acts on estimated
positions---the true minimum separation is dragged to \SI{0.75}{\meter} with
$\sim$\SI{0.8}{\second} per run below the design margin. The DMHE is therefore
not an estimation nicety: it is what keeps the safety layer's world model
accurate enough for its invariance guarantee to bind the physical vehicles.

\section{Computational study of the embedded per-agent solvers}
\label{sec:sim-compute}

The studies above run the reference (Python) coordination layer. This section
measures the \emph{embedded} controller and estimator of \cref{ch:peragent},
instantiated on the drone models of this report through their model interface
and driven in closed loop by a C benchmark harness. Three per-agent
problems are timed: the translational MPC exactly as posed by the coordination
layer ($n_x{=}6$, $n_u{=}3$, $N{=}25$, $12$ constraint rows), the \emph{full
$13$-state quaternion hexacopter NMPC} ($n_u{=}4$, $N{=}20$, $14$ stage and $6$
terminal rows, thrust/torque bounds and velocity state constraints active during
the manoeuvre), and the translational MHE of the DMHE layer ($n_x{=}6$,
$n_w{=}3$ constrained disturbance inputs, $M{=}8$). Timings are medians of
best-of-seven runs on a single core of a shared cloud vCPU
(\texttt{gcc -O3}, double precision, no parallelism); \cref{tab:compute}
reports them.

\begin{table}[t]
\centering
\caption{Per-agent solvers on the drone problems (single core, medians, best of
$7$; ``warm'' is the solvers' warm-started mode, ``exact'' rebuilds the
problem from scratch every sample). The applied-input/instant-estimate latency
is the measurement-to-output path.}
\label{tab:compute}
\small
\begin{tabular}{@{}llccccc@{}}
\toprule
Problem & Mode & Prepare & Feedback & Latency & Inner & Accuracy\\
 & & [\si{\micro\second}] & [\si{\micro\second}] & [\si{\nano\second}] & iter. & \\
\midrule
Translational MPC & exact & 15.3 & 4.7 & 30 & 1.00 & KKT $3.6{\times}10^{-8}$\\
\ ($n_x{=}6$, $N{=}25$) & warm & \textbf{2.7} & 5.5 & 30 & 1.00 & identical\\
\midrule
Quaternion NMPC & exact & 101.6 & 11.0 & 30 & 1.09 & KKT $2.8{\times}10^{-5}$\\
\ ($n_x{=}13$, $N{=}20$) & warm & \textbf{8.5} & 13.2 & 30 & 1.09 & identical\\
\midrule
MHE & exact & 13.7 & 2.5 & 120 & 1.00 & RMSE 0.283 m\\
\ ($n_x{=}6$, $M{=}8$) & warm & \textbf{5.5} & 2.4 & 120 & 1.00 & identical\\
\bottomrule
\end{tabular}
\end{table}

Three observations. First, the absolute numbers: one complete constrained cycle
of the \emph{full nonlinear quaternion-hexacopter} NMPC costs
$\sim$\SI{115}{\micro\second} even when the problem is rebuilt from scratch every
sample, and the applied input is available \SI{30}{\nano\second} after the
state estimate arrives; the per-agent MHE cycle costs
$\sim$\SI{8}{\micro\second}. At the \SI{20}{\hertz} coordination rate used in
this chapter, the entire per-agent estimation-and-control core occupies well
under $0.1\%$ of one embedded core, and even a \SI{1}{\kilo\hertz} attitude-rate
deployment of the full NMPC would use a few percent. Second, the warm-start
behaviour: reusing the work that remains valid from the previous sample cuts the
median preparation by $5.7\times$ (translational), $11.9\times$ (quaternion) and
$2.5\times$ (MHE), while the 99th-percentile preparation stays bounded by the
occasional full refresh. Third, and decisive for a
safety-critical system: the accuracy columns are \emph{identical} across
modes---same KKT residuals, same millimetre-level closed-loop convergence, same
estimation RMSE, and worst-case constraint violation at the
$10^{-4}$ level in normalised units---so the speed-up is not bought with
optimality. Solver failures over all runs: none; the inner solver loop
terminated in at most $4$ passes. It is this combination---microsecond
constrained solves, nanosecond feedback latency, deterministic memory
(\SIrange{32}{209}{\kilo\byte} static per solver) and bounded per-sample
work---that makes distributed constrained estimation and control of the kind
studied here practical on multi-agent UAV hardware. Isolated timings, however,
establish only that the solvers are fast enough; \cref{sec:sim-inloop} closes
the argument by placing them in the loop.

\section{The embedded solvers in the loop}
\label{sec:sim-inloop}

The final hexacopter study replaces the reference-layer estimator and
controller with the \emph{embedded C implementations themselves}, called
per-agent through the interface of \cref{sec:capi} while everything else---the
nonlinear quaternion plants, the measurement generation, the consensus, the CBF
safety filter, the flatness/attitude cascade---remains identical, with the
noise streams paired across configurations. The loop is exactly the deployed
signal path:
\begin{equation*}
\text{nonlinear plant}\ \rightarrow\ \text{embedded MHE}\ \rightarrow\
\text{embedded NMPC}\ \rightarrow\ \text{CBF}\ \rightarrow\ \text{plant}.
\end{equation*}
Validation proceeded in two stages. A single-hexacopter loop was validated
first: over a \SI{20}{\second} take-off-and-track mission on the same noise
stream, the C and Python closed-loop trajectories differ by \SI{0.17}{\meter}
mean (\SI{0.37}{\meter} maximum) over a multi-metre path, with the embedded
estimator's error \emph{half} the reference filter's
(\SI{0.088}{\meter} vs.\ \SI{0.182}{\meter}: the moving window out-smooths the
one-step filter). The full six-agent mission of \cref{sec:sim-hexa} was then run
in four configurations---reference layer; embedded estimator only; embedded
controller only; both---over the same $15$ noise seeds each ($60$ runs).
\Cref{tab:inloop} reports the ablation.

\begin{table}[t]
\centering
\caption{Embedded-solver-in-the-loop equivalence study: the full six-agent
mission (nominal scenario, $3/6$ GNSS-denied for \SI{10}{\second}) in four
configurations over the same $15$ noise seeds. In-loop solver times are medians
of every call made during the missions (prepare${+}$feedback, single shared
core, alongside the Python harness). Divergences and safety-QP failures: zero
in all $60$ runs.}
\label{tab:inloop}
\small
\begin{tabular}{@{}lccccc@{}}
\toprule
Configuration & \shortstack{Denied RMSE\\mean\,/\,med (p95) [m]} &
\shortstack{Shape\\err.\ [m]} & \shortstack{Min sep.\\{[m]}} &
\shortstack{Margin\\breach} & \shortstack{In-loop solver\\median [\si{\micro\second}]}\\
\midrule
Reference layer (baseline) & 0.091\,/\,0.090 (0.100) & 0.092 & $\ge$1.55 & 0 & ---\\
Embedded MHE only & 0.061\,/\,0.061 (0.069) & 0.097 & $\ge$1.54 & 0 & MHE $11.7{+}3.9$\\
Embedded NMPC only & 0.092\,/\,0.091 (0.101) & 0.096 & $\ge$1.53 & 0 & NMPC $47.8{+}9.4$\\
Embedded MHE $+$ NMPC & 0.062\,/\,0.061 (0.070) & 0.096 & $\ge$1.53 & 0 & $\approx76$ total\\
\bottomrule
\end{tabular}
\end{table}

The table reads as an equivalence proof by ablation. Substituting the embedded
\emph{controller} changes nothing statistically---denied-agent RMSE
$0.091\to0.092$\,m, shape error within $0.004$\,m, the same safety margins---%
as expected, since the embedded NMPC solves the same per-agent problem exactly
(and additionally enforces the velocity envelope as hard state constraints).
Substituting the embedded \emph{estimator} leaves formation and safety
untouched while \emph{improving} the estimation error by a third
($0.091\to0.061$\,m: window smoothing over $M{=}8$ intervals against a
recursive one-step filter). The full embedded path combines both effects:
statistically indistinguishable closed-loop behaviour, slightly better
estimation, zero margin breaches, zero divergences---while every estimate and
every command in the loop is produced by the deployed C implementations, whose
in-loop cost totals $\approx$\SI{76}{\micro\second} per agent per cycle with
the measurement-to-command latency still in the tens of nanoseconds. The
multi-hop anchoring extension runs unchanged through the embedded path, since
it only shapes the pseudo-measurement handed to the estimator
($5/6$-denied with the full embedded loop: \SI{0.065}{\meter} denied-agent
RMSE, margin intact). The claim this study licenses is therefore no longer that
the solvers are individually fast enough, but that \emph{the complete
distributed estimation-and-control loop has been validated with the embedded
solver implementations in closed loop on the nonlinear multi-agent plant}.

\section{Fixed-wing cluster}
\label{sec:sim-fw}

\paragraph{Setup.} Four Aerosonde aircraft (\cref{tab:fw-params}) fly a
coordinated loiter. A virtual leader traverses a circular path of radius
\SI{180}{\meter} at \SI{25}{\meter\per\second}; the followers hold a
displacement-based echelon, offset from the leader in along-track phase and in
loiter radius. Each aircraft runs the cascaded autopilot of \cref{ch:formation}---%
airspeed by throttle, altitude through a pitch loop and elevator, course through a
roll loop and aileron, with a coordinated-turn rudder---driven by look-ahead
carrot guidance toward its moving slot, on the full $12$-state nonlinear plant
from the numerically computed trim.

\paragraph{Results.} \Cref{fig:fw-traj} shows the formation: the four aircraft
settle onto concentric loiter circles---the radial echelon---and hold altitude
at \SI{120}{\meter} after the initial transient decays within about
\SI{25}{\second} (transient altitude excursion below \SI{4.2}{\meter};
steady-state altitude error below \SI{0.01}{\meter}).
\Cref{fig:fw-formation} shows the mean slot-tracking error converging to
\SI{9.5}{\meter}---a steady guidance lag of $5\%$ of the loiter radius,
sustained through the continuous coordinated turn. Repeating the run from ten
randomised initial-condition scatters ($\sigma=\SI{3}{\meter}$) converges to the
same coordinated limit cycle in every case (steady-state spread below
\SI{0.01}{\meter}), so the formation is an attractor of the guidance loop rather
than an artefact of one initialisation. The study confirms that the same
coordination architecture drives a fixed-wing cluster---which cannot hover, must
hold airspeed and turns only by banking---through the vehicle-specific autopilot
and guidance, with no change to the distributed logic.

\begin{figure}[t]
\centering
\includegraphics[width=0.92\linewidth]{figs/fig_fw_traj.pdf}
\caption{Fixed-wing cluster on a coordinated loiter: top view (left), showing the
four aircraft on concentric echelon circles, and altitude (right), held at
\SI{120}{\meter}.}
\label{fig:fw-traj}
\end{figure}

\begin{figure}[t]
\centering
\includegraphics[width=0.74\linewidth]{figs/fig_fw_formation.pdf}
\caption{Fixed-wing formation slot-tracking error, converging and held through
the continuous coordinated turn.}
\label{fig:fw-formation}
\end{figure}

\section{Discussion}
\label{sec:sim-discussion}

\Cref{tab:headline} consolidates the campaign. Across both vehicle classes the
framework behaves as the theory predicts, and the Monte-Carlo evidence sharpens
each claim into a statistical one. The consensus layer agrees on the shared
reference within a control period on every tested graph; the model-predictive
controllers track their slots within actuator and tilt limits in all $126$ runs
with no divergence; the CBF safety filter keeps the collision-avoidance set
forward-invariant whenever its world model is accurate, and the campaign
quantifies precisely when it is not---the margin erodes only under broken anchor
coverage or with the estimator ablated, and physical contact never occurs. The
distributed estimator converts relative sensing into resilience with a p95
denied-agent error under \SI{0.13}{\meter} through $5/6$ denial once multi-hop
anchoring restores coverage, and the computational study shows the per-agent
machinery executing constrained nonlinear solves in microseconds with
warm-start factors of $6$--$12\times$ at identical accuracy. The in-the-loop
study closes the gap between those two kinds of evidence: the deployed C
implementations, substituted into the closed loop under paired noise,
reproduce the reference behaviour to within statistical noise---improving the
estimation---so the timing figures and the behavioural figures describe one
and the same computational path. The results are produced on the full nonlinear plants with
noisy measurements, so they exercise the genuine attitude--translation coupling
of the quaternion multirotor model and the stall-limited, minimum-airspeed
behaviour of the fixed-wing model, not linearised surrogates.

\begin{table}[t]
\centering
\caption{Consolidated headline results. RMSE values are Monte-Carlo means
(p95 in parentheses) over the scenario's seeds; compute times are single-core
medians (per-solver breakdown in \cref{tab:compute}).}
\label{tab:headline}
\footnotesize
\setlength{\tabcolsep}{4pt}
\renewcommand{\arraystretch}{1.1}
\begin{tabular}{@{}p{0.30\linewidth}
  >{\centering\arraybackslash}p{0.135\linewidth}
  >{\centering\arraybackslash}p{0.095\linewidth}
  >{\centering\arraybackslash}p{0.085\linewidth}
  >{\centering\arraybackslash}p{0.185\linewidth}@{}}
\toprule
Scenario & Denied RMSE [m] & Min sep.\ [m] & Shape err.\ [m] & Per-agent compute\\
\midrule
Hexacopter nominal, $N_a{=}6$, \SI{10}{s} blackout $3/6$ (30 seeds) &
0.091 (0.099) & $\ge$1.55 & 0.094 &
\multirow{2}{*}{\shortstack[c]{MPC 2.7--15\,\si{\micro\second}\\
MHE 5.5--14\,\si{\micro\second}}}\\[1mm]
\ + denial up to $4/6$ (48 seeds) & $\le$0.106 (0.117) & $\ge$1.49 & $\le$0.101 & \\
\midrule
Anchor coverage broken ($5/6$ denied; ring graph) &
0.45--0.51 (0.69--0.80) & 0.90--0.98 & 0.24--0.27 & unchanged\\
\ + multi-hop anchoring (24 seeds) & 0.11--0.12 (0.13) & $\ge$1.49 & $\le$0.10 & unchanged\\
\midrule
Ablation: controller on dead-reckoning filter &
(0.91 in loop) & 0.75 & 0.449 & unchanged\\
\midrule
Embedded C solvers in the loop (15 paired seeds) &
0.062 (0.070) & $\ge$1.53 & 0.096 &
$\approx$76\,\si{\micro\second}/agent, in loop\\
\midrule
Fixed-wing echelon loiter, $N_a{=}4$ (10 init.\ seeds) &
--- & spacing held & slot lag 9.5 & autopilot layer\\
\bottomrule
\end{tabular}
\end{table}
